% ============ 第2章 §2.4 赋范空间 (原书 pp.51-54) ============

\section{赋范空间}\label{sec:2.4}
\index{0symbolX*@$X^{*}$}

\index{0symbolAnormB@$\lVert A \rVert_{B(X,Y)}$}

\index{0symbolAnormXY@$\lVert A \rVert_{X \to Y}$}

\index{J@距离函数}

\index{J@距离空间}

\index{S@算子范数}

\index{X@线性泛函}

\index{D@对偶空间}

\index{F@赋范空间}

\index{B@Banach 空间}






















从这里开始转入函数解析的内容. 关于 Lebesgue 积分, 读本书时希望读者了解的事项已不作证明地给出, 但函数解析将给出证明. 函数解析是非常广范的数学领域, 书籍极多, 但本书不使用很多函数解析知识, 所以简洁地总结函数解析所需的知识.

\subsection*{距离空间, 赋范空间, Banach 空间}

本书因处理 Fourier 变换的关系, 对复值 Banach 空间感兴趣. 首先复习距离空间的定义. $(X, d)$ 是距离空间是指满足以下条件.

(1) $X$ 是集合.

(2) $d: X \times X \to [0, \infty)$ 是称为距离函数的映射. 即, 对所有 $x, y, z \in X$ 以下条件成立.

(A) $d(x, x) = 0$, 且 $x \neq y$ 时 $d(x, y) > 0$.

(B) $d(x, y) = d(y, x)$.

(C) $d(x, z) \leq d(x, y) + d(y, z)$. 这个不等式称为三角不等式.

作为距离空间的特殊例子, 赋范空间与具备完备性的赋范空间即 Banach 空间重要.

\begin{definition}[赋范空间]\label{def:2.18}
设 $X$ 为复线性空间. 映射 $\|\cdot\|: x \in X \mapsto \|x\| \in [0, \infty)$ 是范数是指满足以下 $2$ 个条件.

(1) 对 $x, y \in X$, 令 $d(x, y) = \|x - y\|$ 时, $d$ 是距离函数.

(2) 对所有 $x \in X$ 与 $a \in \mathbb{C}$, $\|ax\| = |a|\cdot\|x\|$ 成立.
\end{definition}

配备范数的复线性空间称为复赋范空间. 实赋范空间等也通过替换系数体来定义. 不特别注明系数体时简称为赋范空间.

考虑数列空间时没有深入思考完备性, 但无限维赋范空间未必完备 (参照练习问题~\ref{exer:1.30}), 所以明确说具有完备性很重要. 基于这一点, 定义 Banach 空间.

\begin{definition}[Banach 空间]\label{def:2.19}
在赋范空间中, 当对应的距离完备时, 该空间称为 Banach 空间. 即, 赋范空间 $X$ 是 Banach 空间是指满足以下条件: 对任意 $\varepsilon > 0$, 存在某个自然数 $N$, 使得若 $m, n > N$ 则 $\|x_{m} - x_{n}\| < \varepsilon$ 的 $X$ 的点列 $\{x_{n}\}_{n=1}^{\infty}$ 收敛.
\end{definition}

由于属于测度论的范畴, 下面的定理省略证明.

\begin{theorem}\label{thm:2.21}
当 $1 \leq p \leq \infty$, $(X, \mathcal{B}, \mu)$ 为测度空间时, $L^{p}(\mu)$ 是 Banach 空间.
\end{theorem}

\begin{proof}
这里不作证明. 请参考适当的文献.
\end{proof}

\subsection*{有界线性算子}

函数解析中重要的概念有有界线性算子. 关于这个概念, 已通过数列看过种种例子. 调和分析学中, 研究函数与算子的性质, 通过有界线性算子这一概念表达其结论的情形很多.

\begin{definition}\label{def:2.20}
设 $X$, $Y$ 为赋范空间. 线性算子 $A: X \to Y$ 是有界算子, 是指
\[
\|A\|_{X \to Y} = \|A\|_{B(X, Y)} = \sup\{\|Ax\|_{Y} : \|x\|_{X} = 1\}
\]
有限. 这个量称为算子范数. 把有界线性算子 $A: X \to Y$ 的全体构成的集合写作 $B(X, Y)$. $X = Y$ 时省略 $Y$, 写作 $B(X)$. $Y = \mathbb{C}$ 时, 把 $B(X, Y)$ 写作 $X^{*}$, 称为 $X$ 的对偶空间, $X^{*}$ 的元称为线性泛函.
\end{definition}

与赋范空间一样, 对实赋范空间 $X$, 也用相同记号 $X^{*}$ 表示 $B(X, \mathbb{R})$.

研究数列空间时没有意识到这一点, 但下面连接连续性与有界性的定理是基本的.

\begin{theorem}\label{thm:2.22}
设 $X$, $Y$ 为赋范空间. 线性算子 $A: X \to Y$ 是有界算子与 $A: X \to Y$ 连续等价.
\end{theorem}

\begin{proof}
假设 $A$ 有界, 由
\[
\|Ax - Ay\|_{Y} = \|A(x - y)\|_{Y} \leq \|A\|_{X \to Y}\|x - y\|_{X} \quad (x, y \in X)
\]
可知 $A$ 的连续性.

若 $A$ 连续, 则 $\{x \in X : \|Ax\|_{Y} < 1\} = A^{-1}(\{y \in Y : \|y\|_{Y} < 1\})$ 是包含 $0$ 的 $X$ 的开集合. 因此存在某个 $\delta > 0$, 若 $\|x\|_{X} < \delta$ 则 $\|Ax\|_{Y} < 1$. 由此得到 $\|A\|_{X \to Y} \leq \delta^{-1}$.
\end{proof}

用下面的方法构造有界线性算子的情形很多. $D$ 在 $X$ 中稠密是指 $\overline{D} = X$, 即 $D$ 的闭包与 $X$ 一致.

\begin{theorem}\label{thm:2.23}
设 $K > 0$ 为常数, $X$, $Y$ 为 Banach 空间, $D$ 是 $X$ 的稠密子集且为线性空间. $A$ 是从 $D$ 到 $Y$ 的线性算子, 若对所有 $x \in D$ 满足 $\|Ax\|_{Y} \leq K\|x\|_{X}$, 则存在唯一的有界线性算子 $B: X \to Y$, 对所有 $x \in D$ 有 $Bx = Ax$, 且对所有 $x \in X$ 有 $\|Bx\|_{Y} \leq K\|x\|_{X}$ 成立.
\end{theorem}

把上述两处 $\leq K$ 改写为 $= K$ 也可以, 但那种情况下的证明类似, 省略.

\begin{proof}
首先构造有界线性算子 $B: X \to Y$. 对 $x \in X$, 由 $D$ 的稠密性, 存在收敛于 $x$ 的 $D$ 的点列 $\{x_{k}\}_{k=1}^{\infty}$. 若 $k, k' \in \mathbb{N}$, 则 $\|Ax_{k} - Ax_{k'}\|_{Y} = \|A(x_{k} - x_{k'})\|_{Y} \leq K\|x_{k} - x_{k'}\|_{X}$, 所以 $\{Ax_{k}\}_{k=1}^{\infty}$ 在 $Y$ 中收敛. 又, 对收敛于 $x$ 的 $D$ 的其他点列 $\{v_{k}\}_{k=1}^{\infty}$, 若 $k \in \mathbb{N}$, 则 $\|Ax_{k} - Av_{k}\|_{Y} = \|A(x_{k} - v_{k})\|_{Y} \leq K\|x_{k} - v_{k}\|_{X}$, 所以 $\{Av_{k}\}_{k=1}^{\infty}$ 收敛到与 $\{Ax_{k}\}_{k=1}^{\infty}$ 相同的收敛目标. 因此, 与收敛于 $x$ 的 $D$ 的点列 $\{x_{k}\}_{k=1}^{\infty}$ 的取法无关, 可定义 $Bx = \lim_{k \to \infty}Ax_{k}$. 若 $x \in D$, 则对 $k = 1, 2, \ldots$ 令 $x_{k} = x$, $\{x_{k}\}_{k=1}^{\infty}$ 收敛于 $x$, 所以可知 $Bx = Ax$. 又, 对这种 $\{x_{k}\}_{k=1}^{\infty}$,
\[
\|Bx\|_{Y} = \lim_{k \to \infty}\|Ax_{k}\|_{Y} \leq K\lim_{k \to \infty}\|x_{k}\|_{X} = K\|x\|_{X}
\]
成立.

最后, 由 $D$ 的稠密性, 连续映射 $B$ 的定义由上述方法决定, 所以是唯一的.
\end{proof}

$X = Y = L^{p}(\mathbb{R}^{n})$ 时, 从 $X$ 到 $Y$ 的有界线性算子称为 $L^{p}$-有界 (线性) 算子.

关于算子的合成, 下面的公式经常使用.

\begin{theorem}\label{thm:2.24}
设 $X$, $Y$, $Z$ 为 Banach 空间. 对有界线性算子 $A: X \to Y$, $B: Y \to Z$, $\|BA\|_{X \to Z} \leq \|A\|_{X \to Y}\|B\|_{Y \to Z}$ 成立.
\end{theorem}

\begin{proof}
由于 $\|BAx\|_{Z} = \|B(Ax)\|_{Z} \leq \|B\|_{Y \to Z}\|Ax\|_{Y} \leq \|B\|_{Y \to Z}\|A\|_{X \to Y}\|x\|_{X}$, $x \in X$, 所以显然.
\end{proof}

举一个常用的例子.

\begin{example}\label{exa:2.4}
设 $(X, \mathcal{M}, \mu)$ 为 $\sigma$-有限测度空间, $g$ 为 $L^{\infty}(\mu)$ 函数. 此时, 对 $1 \leq p \leq \infty$, 乘法算子 $f \in L^{p}(\mu) \mapsto gf \in L^{p}(\mu)$ 满足 $\|gf\|_{L^{p}(\mu)} \leq \|g\|_{L^{\infty}(\mu)}\|f\|_{L^{p}(\mu)}$, 所以是有界线性算子, 但若 $N < \|g\|_{L^{\infty}(\mu)}$, 则 $E = \{x \in X : |g(x)| > N\}$ 具有正的测度, 特别由 $\sigma$-有限性, 存在可测集合 $F$ 使得 $0 < \mu(E \cap F) < \infty$ 成立. 令 $f = \chi_{F}$ 来考虑即可明白, 这个乘法算子的范数超过 $N$. 因此, 这个乘法算子的范数是 $\|g\|_{L^{\infty}(\mu)}$.
\end{example}
